% Documentación técnica — Sistema domótico de la casa (openHAB + LCN + solar)
% Compilar: pdflatex documentacion-tecnica.tex  (dos veces, para índice y referencias)
\documentclass[11pt,a4paper]{report}

\usepackage[T1]{fontenc}
\usepackage[utf8]{inputenc}
\usepackage{lmodern}
\usepackage[spanish,es-noshorthands,es-tabla]{babel}
\usepackage[margin=2.3cm]{geometry}
\usepackage{amsmath}
\usepackage{xcolor}
\usepackage{graphicx}
\usepackage{float}
\usepackage{booktabs}
\usepackage{longtable}
\usepackage{tabularx}
\usepackage{array}
\usepackage{enumitem}
\usepackage{fancyhdr}
\usepackage{listings}
\usepackage{tikz}
\usetikzlibrary{positioning,shapes.geometric,shapes.misc,shapes.symbols,arrows.meta,fit,backgrounds,calc,shadows}
\usepackage[colorlinks=true,linkcolor=azul,urlcolor=azul]{hyperref}

% ---------------------------------------------------------------- colores
\definecolor{azul}{HTML}{1F5FA8}
\definecolor{solar}{HTML}{E8A317}
\definecolor{verde}{HTML}{2E8B57}
\definecolor{rojo}{HTML}{C0392B}
\definecolor{gris}{HTML}{5D6D7E}
\definecolor{lila}{HTML}{7D3C98}
\definecolor{cian}{HTML}{138D90}
\definecolor{fondo}{HTML}{F4F6F8}

\newcommand{\code}[1]{\texttt{\detokenize{#1}}}

\pagestyle{fancy}
\fancyhf{}
\fancyhead[L]{\small\textcolor{gris}{Sistema domótico — Documentación técnica}}
\fancyhead[R]{\small\textcolor{gris}{\leftmark}}
\fancyfoot[C]{\thepage}
\renewcommand{\headrulewidth}{0.4pt}

\lstset{basicstyle=\ttfamily\small,breaklines=true,frame=single,backgroundcolor=\color{fondo},
  rulecolor=\color{gris!40},columns=fullflexible,keepspaces=true,
  literate={á}{{\'a}}1 {é}{{\'e}}1 {í}{{\'i}}1 {ó}{{\'o}}1 {ú}{{\'u}}1 {ñ}{{\~n}}1}

% ---------------------------------------------------------------- estilos TikZ
\tikzset{
  >={Stealth[length=2.2mm]},
  caja/.style={draw=#1!80!black,fill=#1!12,rounded corners=3pt,thick,align=center,
               minimum height=9mm,inner sep=4pt,font=\small},
  caja/.default=azul,
  nube/.style={draw=#1!80!black,fill=#1!10,cloud,cloud puffs=11,aspect=2.2,thick,align=center,font=\small,inner sep=1pt},
  nube/.default=gris,
  disp/.style={draw=#1!80!black,fill=white,rounded corners=2pt,align=center,font=\scriptsize,
               minimum width=17mm,minimum height=7mm},
  disp/.default=gris,
  grupo/.style={draw=#1!70!black,dashed,rounded corners=6pt,inner sep=8pt},
  grupo/.default=gris,
  etq/.style={font=\scriptsize\itshape,text=gris},
  flujo/.style={->,thick,draw=#1},
  flujo/.default=gris,
  bus/.style={line width=2.2pt,draw=verde},
  estado/.style={draw=#1!80!black,fill=#1!15,circle,minimum size=21mm,thick,align=center,font=\small},
  linea/.style={draw=gris!70,dashed},
  activ/.style={fill=azul!20,draw=azul,minimum width=2.4mm},
}

\title{\vspace{-1cm}\textcolor{azul}{\Huge\bfseries Sistema domótico de la casa}\\[6pt]
       \Large Documentación técnica y de ingeniería de software\\[4pt]
       \large openHAB 5.2.3 · bus LCN · energía solar · cámaras · Tuya · eWeLink}
\author{Ismael Minchala}
\date{Versión 1.1 — 19 de noviembre de 2026}

\begin{document}

\begin{titlepage}
\centering
\vspace*{1.5cm}
{\color{azul}\rule{\linewidth}{1.2pt}}\\[0.6cm]
{\Huge\bfseries\color{azul} Sistema domótico de la casa}\\[0.4cm]
{\Large Documentación técnica y de ingeniería de software}\\[0.3cm]
{\color{azul}\rule{\linewidth}{1.2pt}}\\[1.2cm]

\begin{tikzpicture}[scale=1.0]
  \node[caja=azul,minimum width=34mm,minimum height=14mm,font=\large\bfseries] (oh) at (0,0) {openHAB\\[-2pt]\small 5.2.3};
  \node[caja=verde,minimum width=26mm] (lcn) at (-4.6,1.9) {Bus LCN\\\scriptsize 7 módulos};
  \node[caja=solar,minimum width=26mm] (sol) at (4.6,1.9) {Energía solar\\\scriptsize DessMonitor};
  \node[caja=lila,minimum width=26mm] (cam) at (-4.6,-1.9) {Cámaras\\\scriptsize NVR Amcrest};
  \node[caja=cian,minimum width=26mm] (tuya) at (4.6,-1.9) {Tuya · eWeLink · TV\\\scriptsize Wi-Fi · aspiradora};
  \node[caja=rojo,minimum width=26mm] (app) at (0,3.2) {App iPhone · MainUI};
  \node[caja=gris,minimum width=26mm] (db) at (0,-3.2) {InfluxDB · Grafana};
  \foreach \n/\c in {lcn/verde,sol/solar,cam/lila,tuya/cian,app/rojo,db/gris}
    \draw[<->,very thick,draw=\c] (oh) -- (\n);
\end{tikzpicture}

\vfill
{\large Ismael Minchala}\\[4pt]
{\large Versión 1.1 — 19 de noviembre de 2026}\\[8pt]
{\small Repositorio privado: \url{https://github.com/ismael-minchala/casa-domotica}}
\end{titlepage}

\section*{Historial de versiones}
\begin{tabularx}{\linewidth}{@{}llX@{}}
\toprule
\textbf{Versión} & \textbf{Fecha} & \textbf{Cambios}\\\midrule
1.0 & 3/11/2026 & Primera versión: LCN, energía solar, ahorro automático, cámaras, Tuya y TV.\\
1.1 & 19/11/2026 & Portón y luz exterior (eWeLink), aspiradora Roborock, autonomía de la batería,
  apagado general y monitor de luces, vigilancia de fotos congeladas de las cámaras, detección de datos
  congelados de DessMonitor, monitor de red; inventario, red, reglas e incidentes actualizados.\\
\bottomrule
\end{tabularx}

\tableofcontents

% =====================================================================
\chapter{Introducción}

\section{Propósito}
Este documento describe la arquitectura, la topología de red, los componentes de software, los flujos de
datos, la operación y el mantenimiento del sistema domótico de la casa. Sirve como referencia para
operar el sistema, diagnosticar fallos, restaurarlo desde el respaldo y planificar su evolución.

\section{Alcance}
El sistema integra en una sola plataforma (openHAB) los siguientes subsistemas:
\begin{itemize}[nosep]
  \item \textbf{Iluminación y persianas} sobre el bus cableado \textbf{LCN} (7 módulos, pasarela LCN-PKE).
  \item \textbf{Energía solar}: inversor aislado de 24\,V con baterías y controlador de carga, monitorizado
        mediante la API en la nube de DessMonitor; historial en InfluxDB y paneles en Grafana.
  \item \textbf{Ahorro automático de energía}: reglas que atenúan o apagan cargas según el estado de la batería,
        y estimación de la autonomía restante de la batería.
  \item \textbf{Videovigilancia}: grabador NVR Amcrest con 3 cámaras activas, con reinicio automático si la
        imagen se congela.
  \item \textbf{Dispositivos Wi-Fi Tuya}: interruptores, atenuador del garaje, panel y sensor de gas.
  \item \textbf{Dispositivos Sonoff (nube eWeLink)}: apertura y cierre del portón y luz exterior.
  \item \textbf{Aspiradora robot} Roborock S5 (Xiaomi, binding \code{miio}).
  \item \textbf{Salir de casa}: apagado general de luces y televisores y monitor de luces encendidas.
  \item \textbf{Multimedia}: Google TV / Android TV y Chromecast.
  \item \textbf{Monitor de la red} doméstica (latencia, DNS, velocidad y estado del servidor) con tablero Grafana.
  \item \textbf{Acceso móvil} (app openHAB para iOS, local y remoto vía myopenHAB) y notificaciones push.
\end{itemize}

\section{Convenciones}
\begin{itemize}[nosep]
  \item Las direcciones IP sin prefijo pertenecen a la red \code{192.168.18.0/24} (p.\,ej. \code{.30} = \code{192.168.18.30}).
  \item Los nombres de \emph{Items}, \emph{Things} y archivos se escriben en \code{monoespaciado}.
  \item Este documento \textbf{no contiene credenciales}. En el repositorio aparecen como \code{<REDACTED>}.
\end{itemize}

\section{Glosario}
\begin{longtable}{@{}p{3.2cm}p{11.8cm}@{}}
\toprule
\textbf{Término} & \textbf{Definición}\\\midrule\endhead
Binding & Complemento de openHAB que habla el protocolo de una tecnología (LCN, Tuya, ipcamera\ldots).\\
Thing & Representación de un dispositivo físico o servicio dentro de openHAB.\\
Channel & Función concreta de un Thing (salida de un módulo, volumen de un TV\ldots).\\
Item & Variable de openHAB enlazada a un Channel; es lo que ven las interfaces y las reglas.\\
Sitemap & Descripción textual de la interfaz para la app móvil nativa.\\
MainUI & Interfaz web de openHAB (páginas, widgets, administración).\\
LCN & \emph{Local Control Network}: bus domótico cableado de Issendorff.\\
PCK & Protocolo de texto de LCN sobre TCP (puerto 4114).\\
PKE / PCHK & Pasarela Ethernet–LCN (hardware LCN-PKE que ejecuta el servicio PCHK 3.2.2).\\
SOC & \emph{State of Charge}: porcentaje de carga de la batería.\\
NVR & Grabador de vídeo en red; concentra las cámaras.\\
MJPEG & Vídeo como secuencia de imágenes JPEG; compatible con la app y los navegadores.\\
Flux & Lenguaje de consultas de InfluxDB 2.\\
eWeLink & Nube de los equipos Sonoff (servidores \code{*-apia.coolkit.cc}).\\
OAuth 2.0 & Protocolo de autorización: el usuario inicia sesión en la web del proveedor y openHAB recibe
            un código que canjea por un token, sin conocer la contraseña.\\
miio & Protocolo local de los equipos Xiaomi (UDP 54321); la clave del equipo se obtiene de la cuenta Xiaomi.\\
Autonomía & Tiempo estimado hasta que la batería llega al mínimo en que el inversor pasa a la red.\\
\bottomrule
\end{longtable}

% =====================================================================
\chapter{Requisitos}

\section{Requisitos funcionales}
\begin{longtable}{@{}p{1.3cm}p{10.2cm}p{3.3cm}@{}}
\toprule
\textbf{ID} & \textbf{Requisito} & \textbf{Estado}\\\midrule\endhead
RF-01 & Encender, apagar y regular cada luz LCN desde la app (interruptor + deslizador de intensidad). & Implementado\\
RF-02 & Subir y bajar la persiana del dormitorio principal. & Implementado\\
RF-03 & Controlar las luces Tuya de los dormitorios, el panel y el garaje. & Implementado (Alejandro sin conexión)\\
RF-04 & Mostrar en tiempo real generación solar, consumo, SOC, modo y alarmas del inversor. & Implementado\\
RF-05 & Diagrama animado del flujo de energía (solar, red, inversor, batería, casa). & Implementado\\
RF-06 & Gráficos históricos y energía diaria generada/consumida (kWh). & Implementado\\
RF-07 & Ahorro automático: atenuar luces al 40\,\% con batería baja y apagar cargas con batería crítica. & Implementado\\
RF-08 & Avisos push al móvil ante cambios de ahorro, alarmas del inversor y pérdida de datos. & Implementado\\
RF-09 & Ver en vivo las cámaras Calle, Entrada e Interior desde la app. & Implementado\\
RF-10 & Controlar TV de Sala, MIBOX3 y TV Alejandro (encendido, volumen, reproducción). & Parcial (falta emparejar PIN)\\
RF-11 & Sincronizar la hora del bus LCN. & Implementado (PKE)\\
RF-12 & Respaldo semanal automático de la configuración, sin secretos. & Implementado\\
RF-13 & Estimar las horas de autonomía de la batería con el consumo de los últimos 30 min. & Implementado\\
RF-14 & Apagar con un botón todas las luces y televisores al salir; ver qué luces siguen encendidas o sin conexión. & Implementado\\
RF-15 & Abrir y cerrar el portón y controlar la luz exterior (Sonoff vía eWeLink). & Implementado\\
RF-16 & Controlar la aspiradora (limpiar, pausar, volver a la base, potencia) y ver mapa y consumibles. & Implementado\\
RF-17 & Medir la calidad de la red doméstica y del enlace a Internet y mostrarla en Grafana. & Implementado\\
\bottomrule
\end{longtable}

\section{Requisitos no funcionales}
\begin{longtable}{@{}p{1.3cm}p{4cm}p{9.5cm}@{}}
\toprule
\textbf{ID} & \textbf{Atributo} & \textbf{Criterio}\\\midrule\endhead
RNF-01 & Latencia & Un comando de luz debe ejecutarse en menos de 1\,s en condiciones normales.\\
RNF-02 & Disponibilidad & openHAB arranca como servicio de Windows; InfluxDB y Grafana arrancan solos.\\
RNF-03 & Robustez & Un flujo de cámara caído no debe bloquear el servidor (Thing deshabilitado $\Rightarrow$ 404 inmediato).\\
RNF-04 & Licencias & No abrir conexiones extra a la PKE (licencia limitada a 3 sesiones PCK).\\
RNF-05 & Seguridad & Las credenciales nunca se escriben en código, documentos ni en el repositorio.\\
RNF-06 & Usabilidad & Interfaz en español, agrupada por plantas y habitaciones.\\
RNF-07 & Mantenibilidad & Configuración por archivos versionados en Git; respaldo y restauración documentados.\\
RNF-08 & Rendimiento & Las cámaras se transcodifican a 1280\,px y 4\,fps ($\approx$\,41\,\% CPU con 3 flujos abiertos).\\
RNF-09 & Autorrecuperación & Una cámara con la foto congelada se reinicia sola (3 min sin cambios, máximo un reinicio cada 15 min).\\
RNF-10 & Veracidad de datos & Una lectura vieja del inversor (más de 15 min) no se publica como nueva.\\
\bottomrule
\end{longtable}

% =====================================================================
\chapter{Arquitectura del sistema}

\section{Diagrama de contexto}
La figura~\ref{fig:contexto} muestra el sistema como una caja negra y los actores externos con los que
interactúa.

\begin{figure}[H]
\centering
\begin{tikzpicture}[node distance=1.2cm and 1.4cm]
  \node[caja=azul,minimum width=48mm,minimum height=22mm,font=\bfseries] (sys) {Sistema domótico\\[-1pt]\normalfont\small servidor openHAB};
  \node[caja=rojo,above left=1.4cm and 0.6cm of sys] (usr) {Habitantes\\\scriptsize app iOS · navegador};
  \node[caja=rojo,above right=1.4cm and 0.6cm of sys] (adm) {Administrador\\\scriptsize MainUI · LCN-PRO};
  \node[nube=gris,left=1.6cm of sys] (cloud) {myopenHAB\\\scriptsize acceso remoto · push};
  \node[nube=solar,right=1.6cm of sys] (dess) {DessMonitor\\\scriptsize API del inversor};
  \node[nube=cian,below right=1.3cm and 0.3cm of sys] (tuyac) {Tuya Cloud\\\scriptsize alta de equipos};
  \node[caja=verde,below left=1.4cm and 0.2cm of sys] (fis) {Equipos físicos\\\scriptsize LCN · cámaras · Tuya · TV};
  \node[nube=gris,below=1.9cm of sys] (gh) {GitHub\\\scriptsize respaldo privado};
  \node[nube=cian,above=2.0cm of sys] (ewl) {eWeLink · Xiaomi\\\scriptsize portón · aspiradora};
  \draw[flujo=cian,<->] (sys) -- node[etq,right]{HTTPS / OAuth} (ewl);
  \draw[flujo=rojo,<->] (usr) -- node[etq,left]{comandos / estado} (sys);
  \draw[flujo=rojo,<->] (adm) -- node[etq,right]{configuración} (sys);
  \draw[flujo,<->] (cloud) -- node[etq,above]{HTTPS} (sys);
  \draw[flujo=solar,<-] (sys) -- node[etq,above]{cada 2 min} (dess);
  \draw[flujo=cian,<->] (sys) -- (tuyac);
  \draw[flujo=verde,<->] (sys) -- node[etq,left]{LAN} (fis);
  \draw[flujo,->] (sys) -- node[etq,right]{semanal} (gh);
\end{tikzpicture}
\caption{Diagrama de contexto.}
\label{fig:contexto}
\end{figure}

\section{Arquitectura lógica en capas}
openHAB separa el acceso al hardware (bindings y Things) de la lógica (Items y reglas) y de la
presentación (sitemap, MainUI). La figura~\ref{fig:capas} resume las capas y sus componentes.

\begin{figure}[H]
\centering
\begin{tikzpicture}[x=1cm,y=1cm]
  \def\capa#1#2#3#4{%
    \fill[#3!10,draw=#3!70!black,rounded corners=4pt,thick] (0,#1) rectangle (15,#1+1.45);
    \node[anchor=west,font=\bfseries\small,text=#3!70!black] at (0.2,#1+1.1) {#2};
    \node[anchor=west,font=\scriptsize,text width=14.4cm] at (0.2,#1+0.45) {#4};}
  \capa{7.6}{Presentación}{rojo}{Sitemap \code{casa} (app iOS) · MainUI: \code{home}, \code{overview}, \code{energia}, \code{energia_graficos}, \code{seguridad}, \code{variables} · HTML en \code{/static}: \code{flujo-energia}, \code{energia-diaria}, \code{ewelink} · Grafana «Energía» y «Red»}
  \capa{5.9}{Lógica y automatización}{lila}{12 reglas JS: \code{dessmonitor}, \code{ahorro-energia}, \code{bateria-autonomia}, \code{casa-general}, \code{camaras-vigilancia}, \code{ewelink} · Avisos openHAB Cloud · \code{monitor-red.ps1} (fuera de openHAB)}
  \capa{4.2}{Modelo de datos}{azul}{131 Items (112 de valor + 19 grupos) · modelo semántico por plantas y habitaciones · grupos \code{gInversor}, \code{gAhorro}, \code{gAspiradora}, \code{gEwelink}\ldots{} · 51 enlaces Item--Channel · metadatos}
  \capa{2.5}{Integración}{verde}{26 Things · bindings \code{lcn}, \code{tuya}, \code{ipcamera}, \code{androidtv}, \code{chromecast}, \code{miio}, \code{sonoff} (solo cuenta), \code{amazonechocontrol} · HTTP a DessMonitor y eWeLink · \code{ffmpeg}}
  \capa{0.8}{Persistencia e infraestructura}{gris}{InfluxDB 2.7 (org \code{Home}, bucket \code{data}) · Windows 10 · Java Zulu 21 · Karaf 4.4 · servicio \code{openHAB} · Git + tareas programadas de respaldo y monitor de red}
  \foreach \y in {2.25,3.95,5.65,7.35} \draw[<->,thick,gris] (7.5,\y) -- ++(0,0.3);
\end{tikzpicture}
\caption{Arquitectura lógica en capas.}
\label{fig:capas}
\end{figure}

\section{Diagrama de componentes}

\begin{figure}[H]
\centering
\resizebox{\linewidth}{!}{%
\begin{tikzpicture}[node distance=0.55cm and 0.6cm]
  % núcleo
  \node[caja=azul,minimum width=36mm,minimum height=12mm] (bus) {Bus de eventos\\\scriptsize openHAB core};
  \node[caja=azul,above=0.7cm of bus,minimum width=36mm] (items) {Registro de Items};
  \node[caja=lila,left=1.6cm of items] (rules) {Motor de reglas\\\scriptsize JS Scripting};
  \node[caja=gris,right=1.6cm of items] (pers) {Persistencia\\\scriptsize influxdb};
  \node[caja=rojo,above=0.8cm of items,minimum width=36mm] (ui) {REST / SSE\\\scriptsize :8080};
  \node[caja=rojo,left=1.1cm of ui] (app) {App iOS\\\scriptsize sitemap};
  \node[caja=rojo,right=1.1cm of ui] (main) {MainUI\\\scriptsize navegador};
  \node[caja=gris,right=1.1cm of pers] (influx) {InfluxDB\\\scriptsize 127.0.0.1:8086};
  \node[caja=gris,above=0.8cm of influx] (graf) {Grafana\\\scriptsize :3000};
  % bindings
  \node[caja=verde,below left=1.1cm and 2.2cm of bus] (b1) {LCN};
  \node[caja=cian,right=0.4cm of b1] (b2) {Tuya};
  \node[caja=lila,right=0.4cm of b2] (b3) {ipcamera};
  \node[caja=cian,right=0.4cm of b3] (b4) {Android TV};
  \node[caja=cian,right=0.4cm of b4] (b5) {Chromecast};
  \node[caja=solar,left=0.5cm of b1] (b0) {HTTP\\\scriptsize DessMonitor};
  \node[caja=cian,right=0.4cm of b5] (b6) {miio\\\scriptsize aspiradora};
  \node[caja=cian,right=0.4cm of b6] (b7) {HTTP\\\scriptsize eWeLink};
  \node[grupo=verde,fit=(b0)(b7),label={[etq]below:bindings y clientes HTTP de las reglas}] {};
  \draw[flujo=azul,<->] (items) -- (bus);
  \draw[flujo=lila,<->] (rules) -- (items);
  \draw[flujo=gris,->] (items) -- (pers);
  \draw[flujo=gris,->] (pers) -- (influx);
  \draw[flujo=gris,<-] (graf) -- node[etq,right]{Flux} (influx);
  \draw[flujo=rojo,<->] (ui) -- (items);
  \draw[flujo=rojo,<->] (app) -- (ui);
  \draw[flujo=rojo,<->] (main) -- (ui);
  \foreach \b in {b0,b1,b2,b3,b4,b5,b6,b7} \draw[flujo=verde,<->] (\b.north) -- (bus.south);
\end{tikzpicture}}
\caption{Componentes principales y sus dependencias.}
\label{fig:componentes}
\end{figure}

\section{Diagrama de despliegue}

\begin{figure}[H]
\centering
\begin{tikzpicture}[node distance=0.4cm]
  \node[caja=azul,minimum width=38mm] (oh) {Servicio \code{openHAB}\\\scriptsize Java Zulu 21 · Karaf · :8080/:8443};
  \node[caja=gris,right=0.5cm of oh,minimum width=34mm] (ifx) {InfluxDB 2.7\\\scriptsize tarea programada (SYSTEM)};
  \node[caja=gris,right=0.5cm of ifx,minimum width=30mm] (gf) {Grafana 13\\\scriptsize servicio · :3000};
  \node[caja=lila,below=0.5cm of oh,minimum width=38mm] (ff) {ffmpeg\\\scriptsize \texttt{C:\textbackslash ffmpeg\textbackslash bin}};
  \node[caja=verde,right=0.5cm of ff,minimum width=34mm] (lpro) {LCN-PRO 6.9\\\scriptsize proyecto \code{Casa.lp3}};
  \node[caja=gris,right=0.5cm of lpro,minimum width=30mm] (git) {Git 2.56\\\scriptsize respaldo semanal};
  \node[grupo=azul,fit=(oh)(gf)(ff)(git),inner sep=12pt,label={[font=\bfseries\small,text=azul]above:PC servidor «CASA-ISMAEL» · Windows 10 Pro · Wi-Fi .254}] (pc) {};
  \node[caja=verde,below=1.6cm of ff,minimum width=34mm] (pke) {LCN-PKE (.30)\\\scriptsize PCHK 3.2.2 · TCP 4114};
  \node[caja=lila,right=0.9cm of pke,minimum width=34mm] (nvr) {NVR Amcrest (.31)\\\scriptsize RTSP 554 · HTTP 80};
  \node[caja=cian,right=0.9cm of nvr,minimum width=30mm] (tuy) {Equipos Tuya\\\scriptsize TCP 6668 · UDP 6666/7};
  \draw[flujo=verde,<->] (pke) -- node[etq,left]{PCK} (pke |- pc.south);
  \draw[flujo=lila,<->] (nvr) -- node[etq,left]{RTSP} (nvr |- pc.south);
  \draw[flujo=cian,<->] (tuy) -- node[etq,left]{Tuya 3.3} (tuy |- pc.south);
\end{tikzpicture}
\caption{Despliegue físico del software.}
\label{fig:despliegue}
\end{figure}

% =====================================================================
\chapter{Topología de red}

\section{Diagrama de red}

\begin{figure}[H]
\centering
\resizebox{\linewidth}{!}{%
\begin{tikzpicture}
  % troncal
  \node[nube=gris,minimum width=26mm] (inet) at (0,0.2) {Internet};
  \node[caja=gris] (rt) at (0,-1.5) {Router Huawei\\\scriptsize .1 · DHCP};
  \node[caja=azul] (srv) at (-5.6,-1.5) {Servidor openHAB\\\scriptsize .254 (Wi-Fi, IP fija)};
  \node[disp] (ip) at (-5.6,0.2) {iPhone\\\scriptsize app openHAB};
  \node[caja=gris] (rep) at (6.1,-1.5) {Repetidor Wi-Fi\\\scriptsize proxy ARP};
  \node[caja=verde] (pke) at (-2.9,-3.3) {LCN-PKE\\\scriptsize .30 · Ethernet};
  \node[caja=lila] (nvr) at (1.7,-3.3) {NVR Amcrest\\\scriptsize .31 · NV2104E-HS};
  \draw[thick] (inet) -- (rt);
  \draw[thick,dashed] (rt) -- node[etq,above]{Wi-Fi} (srv);
  \draw[thick,dashed] (rt) -- node[etq,above]{Wi-Fi} (rep);
  \draw[thick] (rt.south) |- ++(0,-0.45) -| (pke.north);
  \draw[thick] (rt.south) |- ++(0,-0.45) -| (nvr.north);
  \draw[rojo,dashed] (srv) -- (ip);
  % bus LCN
  \draw[bus] (-7.8,-4.6) -- (-1.4,-4.6);
  \draw[bus] (pke.south) -- (pke.south |- 0,-4.6);
  \node[etq,anchor=south west] at (-7.8,-4.6) {bus LCN};
  \foreach \i/\m in {0/M5,1/M6,2/M7,3/M8,4/M9,5/M10,6/M102} {
    \node[disp=verde,minimum width=7mm,inner sep=2pt] (m\i) at (-7.5+1.0*\i,-5.35) {\m};
    \draw[verde,thick] (m\i.north) -- ++(0,0.43);
  }
  \node[etq,text width=7.4cm,align=center] at (-4.5,-6.05) {entrada · cocina · comedor · persianas · sala · alarma · dormitorio};
  % cámaras
  \node[disp=lila] (c1) at (0.5,-5.2) {Interior\\\scriptsize PoE 10.1.1.129};
  \node[disp=lila] (c4) at (2.6,-5.2) {Calle\\\scriptsize PoE 10.1.1.128};
  \node[disp=lila] (c3) at (4.7,-5.2) {Entrada\\\scriptsize Wi-Fi .204};
  \draw[lila,thick] (nvr) -- node[etq,left]{ch1} (c1);
  \draw[lila,thick] (nvr) -- node[etq,right]{ch4} (c4);
  \draw[lila,thick,dashed] (nvr) -- node[etq,right]{ch3} (c3);
  % clientes Wi-Fi
  \foreach \n/\t/\x in {t1/{Panel\\\scriptsize .104}/4.6,t2/{Garaje\\\scriptsize .105}/6.1,t3/{Julián\\\scriptsize .199}/7.6,t4/{Roborock\\\scriptsize .142}/9.1}
    { \node[disp=cian] (\n) at (\x,0.2) {\t}; \draw[cian,dashed] (rep) -- (\n); }
  \foreach \n/\t/\x in {v1/{MIBOX3\\\scriptsize .188}/4.6,v2/{TV Ale\\\scriptsize .166}/6.1,v3/{TV Sala\\\scriptsize .155}/7.6}
    { \node[disp=cian] (\n) at (\x,-3.3) {\t}; \draw[cian,dashed] (rep) -- (\n); }
  \node[disp=cian] (sf) at (9.6,-3.3) {Sonoff\\\scriptsize portón · luz};
  \draw[cian,dashed] (rep) -- (sf);
  \node[etq,text width=3.2cm,align=center] at (9.6,-4.25) {IP por DHCP; control solo por la nube eWeLink};
\end{tikzpicture}}
\caption{Topología física y lógica de la red doméstica.}
\label{fig:red}
\end{figure}

\section{Plan de direccionamiento}
\begin{longtable}{@{}lllp{6.2cm}@{}}
\toprule
\textbf{IP} & \textbf{Equipo} & \textbf{Enlace} & \textbf{Observaciones}\\\midrule\endhead
.1 & Router Huawei & — & Puerta de enlace y DHCP.\\
.30 & LCN-PKE & Ethernet & IP fija; PCHK 3.2.2, firmware 2.06; licencia de 3 conexiones.\\
.31 & NVR Amcrest NV2104E-HS & Ethernet & IP fija; PoE para las cámaras 10.1.1.x.\\
.104 & Tuya «Panel» & Wi-Fi & Interruptor (\code{switch_1}).\\
.105 & Tuya «Garaje» & Wi-Fi & Atenuador (\code{bright_value}).\\
.136 & Tuya «Interruptor ale» & Wi-Fi & Sin conexión desde el 2/11/2026.\\
.142 & Aspiradora Roborock S5 & Wi-Fi & Binding \code{miio} (control local).\\
.155 & TV de Sala (Google TV) & Wi-Fi & Pendiente de emparejar (PIN).\\
.166 & TV Alejandro & Wi-Fi & Chromecast activo; Google TV pendiente de PIN.\\
.188 & MIBOX3 & Wi-Fi & Chromecast activo; Google TV pendiente de PIN.\\
.199 & Tuya «Interruptor Julián» & Wi-Fi & A través del repetidor.\\
.204 & Cámara Entrada ASH42-B & Wi-Fi & IP fija; canal 3 del NVR.\\
.254 & Servidor openHAB & Wi-Fi & IP fija; perfil de red «Público» (pendiente).\\
DHCP & Sonoff (portón, luz ext.) & Wi-Fi & DUAL R2 y 4CH; openHAB no les habla en local, solo vía la nube eWeLink.\\
\bottomrule
\end{longtable}

\section{Puertos y servicios}
\begin{longtable}{@{}lllp{6.5cm}@{}}
\toprule
\textbf{Puerto} & \textbf{Equipo} & \textbf{Protocolo} & \textbf{Uso}\\\midrule\endhead
8080 / 8443 & Servidor & HTTP(S) & openHAB: REST, MainUI, app, \code{/static}, \code{/ipcamera}.\\
8086 & Servidor (local) & HTTP & InfluxDB 2.7, solo \code{127.0.0.1}.\\
3000 & Servidor & HTTP & Grafana.\\
4114 & PKE & TCP (PCK) & Pasarela LCN; openHAB mantiene una sesión permanente.\\
80 & PKE & HTTP & Configuración web de la PKE.\\
554 & NVR & RTSP & Flujos de vídeo por canal (\code{subtype=0} principal, \code{1} secundario).\\
80 & NVR & HTTP/ONVIF & API Amcrest y ONVIF.\\
6668 & Tuya & TCP & Control local cifrado (protocolo 3.3).\\
6666/6667 & Tuya & UDP & Anuncios de descubrimiento de IP.\\
6466/6467 & Google TV & TLS & Control remoto / emparejamiento con PIN.\\
8009 & Chromecast & TLS & Cast.\\
54321 & Roborock & UDP (miio) & Control local cifrado de la aspiradora.\\
443 (saliente) & Nubes & HTTPS & DessMonitor, eWeLink (\code{*-apia.coolkit.cc}), Xiaomi, myopenHAB, pruebas de velocidad.\\
\bottomrule
\end{longtable}

% =====================================================================
\chapter{Subsistemas}

\section{Iluminación y persianas (LCN)}
El bus LCN es la columna vertebral de la iluminación. openHAB se conecta a la pasarela LCN-PKE
(\code{lcn:pckGateway:LCN-PCHK}, modo \code{native200}, \code{timeoutMs=3500}) y expone cada salida
como un Item de tipo \code{Dimmer}.

{\small\begin{longtable}{@{}lp{2.6cm}p{1.6cm}p{8.4cm}@{}}
\toprule
\textbf{Módulo} & \textbf{Ubicación} & \textbf{Salida} & \textbf{Item}\\\midrule\endhead
5 & Entrada & 1 / 2 & \code{Module_5_ENTRADA_Output_1/2} (escalera / laterales)\\
6 & Cocina & 1 / 2 & \code{Module_6_COCINA_Output_1/2} (leds / luz)\\
7 & Comedor & 1 / 2 & \code{Module_7_COMEDOR_Output_1/2}\\
8 & Persianas & 1 / 2 & \code{Module_8_PERSIANAS_Output_1/2} (bajar / subir)\\
9 & Sala & 1 / 2 & \code{Module_9_SALA_Output_1/2}\\
10 & Alarma & 1 & \code{Module_10_ALARMA_Output_1} (sirena)\\
102 & Dormitorio principal & 1 / 2, var.\,4 & \code{Module_102_CUARTO_PRINCIPAL_Output_1/2}, \code{Module102Var4}\\
\bottomrule
\end{longtable}}

\paragraph{Secuencia de un comando.} La figura~\ref{fig:seq-luz} muestra qué ocurre al pulsar el
interruptor de una luz en la app. El interruptor envía \code{ON}/\code{OFF} al mismo Item \code{Dimmer}
que usa el deslizador de intensidad.

\begin{figure}[H]
\centering
\begin{tikzpicture}[x=2.55cm,y=-0.62cm]
  \foreach \i/\n/\c in {0/App iOS/rojo,1/openHAB REST/azul,2/Binding LCN/verde,3/LCN-PKE/verde,4/Módulo LCN/verde} {
    \node[caja=\c,minimum width=22mm] (h\i) at (\i,0) {\n};
    \draw[linea] (\i,0.8) -- (\i,11.2);
  }
  \draw[flujo=rojo] (0,2) -- node[etq,above]{POST /items/X \code{OFF}} (1,2);
  \draw[flujo=azul] (1,3) -- node[etq,above]{ItemCommandEvent} (2,3);
  \draw[flujo=verde] (2,4) -- node[etq,above]{PCK \code{>M000005.A1DI000000}} (3,4);
  \draw[flujo=verde] (3,5) -- node[etq,above]{telegrama bus} (4,5);
  \draw[flujo=verde,dashed] (4,6.3) -- node[etq,above]{acuse} (3,6.3);
  \draw[flujo=verde,dashed] (3,7.3) -- node[etq,above]{\code{-M000005!}} (2,7.3);
  \draw[flujo=verde,dashed] (4,8.5) -- node[etq,above]{estado salida 0\,\%} (2,8.5);
  \draw[flujo=azul,dashed] (2,9.5) -- node[etq,above]{ItemStateChangedEvent} (1,9.5);
  \draw[flujo=rojo,dashed] (1,10.5) -- node[etq,above]{SSE: estado \code{0}} (0,10.5);
  \node[etq,anchor=north,text width=10cm,align=center] at (2,11.5) {si el acuse no llega en 3,5\,s, el binding reintenta el comando};
\end{tikzpicture}
\caption{Diagrama de secuencia: apagar una luz LCN desde la app.}
\label{fig:seq-luz}
\end{figure}

\paragraph{Hora del bus.} Desde el 2/11/2026 la PKE tiene su reloj ajustado (zona
\code{America/Guayaquil}) y la opción «sync LCN bus time» activada, por lo que envía la hora al
bus periódicamente. La PKE no tiene puerta de enlace ni DNS configurados (no usa NTP): tras un corte de
energía hay que verificar su reloj en \url{http://192.168.18.30}.

\section{Energía solar}

\subsection{Hardware}
Inversor aislado de 24\,V (dos baterías) con controlador de carga fotovoltaico separado y potencia
nominal de 3000\,W. Un colector Wi-Fi Eybond envía los datos a la nube DessMonitor.

\subsection{Adquisición de datos}
El inversor no expone vatios; la regla \code{dessmonitor.js} los deriva:
\[
P_\text{solar} = V_\text{batería}\cdot I_\text{carga},\qquad
P_\text{casa} = \frac{\text{Carga (\%)}}{100}\cdot 3000\,\text{W}.
\]
$P_\text{casa}$ solo es válida en «Inverter mode»; con red («Mains mode») el inversor
informa carga 0\,\%.

\begin{figure}[H]
\centering
\resizebox{\linewidth}{!}{%
\begin{tikzpicture}[node distance=0.9cm and 1.25cm]
  \node[caja=solar] (inv) {Inversor\\\scriptsize + controlador};
  \node[caja=solar,right=of inv] (col) {Colector\\\scriptsize Eybond Wi-Fi};
  \node[nube=solar,right=of col] (cloud) {API\\DessMonitor};
  \node[caja=lila,right=of cloud] (js) {\code{dessmonitor.js}\\\scriptsize cron 2 min · SHA-1};
  \node[caja=azul,below=1.1cm of js] (it) {Items \code{gInversor}\\\scriptsize SOC, potencias, modo\ldots};
  \node[caja=gris,left=1.4cm of it] (ifx) {InfluxDB\\\scriptsize bucket \code{data}};
  \node[caja=gris,left=of ifx] (gf) {Grafana\\\scriptsize kWh/día};
  \node[caja=rojo,below=1.0cm of it] (ui) {Flujo animado\\\scriptsize \code{flujo-energia.html}};
  \node[caja=rojo,left=1.4cm of ui] (dd) {Energía diaria\\\scriptsize \code{energia-diaria.html}};
  \node[caja=lila,right=1.2cm of it] (ah) {\code{ahorro-energia.js}};
  \draw[flujo=solar] (inv) -- node[etq,above]{RS232} (col);
  \draw[flujo=solar] (col) -- node[etq,above]{TCP} (cloud);
  \draw[flujo=solar] (js) -- node[etq,above]{HTTPS} (cloud);
  \draw[flujo=solar,dashed] (cloud.south east) to[bend right=15] (js.south west);
  \draw[flujo=azul] (js) -- (it);
  \draw[flujo=gris] (it) -- node[etq,above]{everyChange} (ifx);
  \draw[flujo=gris] (ifx) -- node[etq,above]{Flux} (gf);
  \draw[flujo=rojo] (it) -- node[etq,right]{REST 15 s} (ui);
  \draw[flujo=rojo] (ifx) -- node[etq,left]{historial} (dd);
  \draw[flujo=lila] (it) -- (ah);
\end{tikzpicture}}
\caption{Flujo de datos de energía.}
\label{fig:energia}
\end{figure}

\subsection{Energía diaria}
\code{energia-diaria.html} consulta el historial de \code{Solar_Potencia} y \code{Casa_Consumo} y
calcula la integral por día local (UTC$-5$) con la regla del trapecio, partiendo en medianoche y
descartando huecos mayores de 30 minutos:
\[
E_\text{día} = \sum_{k} \frac{P_k + P_{k+1}}{2}\,(t_{k+1}-t_k), \qquad t_{k+1}-t_k \le 30\,\text{min}.
\]
El resultado coincide con el panel «Energía por día» de Grafana (p.\,ej. 1/11/2026: 5,67\,kWh
generados, 4,17\,kWh consumidos).

\subsection{Detección de datos congelados}
Cuando el colector Eybond deja de subir datos, la nube DessMonitor sigue respondiendo con la última
lectura. Desde el 15/11/2026 \code{dessmonitor.js} usa la marca de tiempo del propio inversor (punto
\code{Timestamp}): si tiene más de 15 minutos, no actualiza los Items de potencia ni SOC, deja en
\code{Energia_UltimaLectura} la hora real de esa lectura y pone en \code{Energia_ApiEstado} «Sin datos
nuevos desde las HH:MM (el datalogger no envía a la nube)». Así los gráficos muestran un hueco en vez de
una línea plana falsa, y la regla de vigilancia del ahorro detecta la pérdida de datos.

\subsection{Autonomía de la batería}
La regla \code{bateria-autonomia.js} (cada 2 minutos y ante cambios de SOC o modo) estima cuánto
tiempo queda hasta que la batería llega al mínimo en que el inversor pasa a la red:
\[
t_\text{autonomía} = \frac{(\text{SOC} - \text{SOC}_\text{mín})\cdot E_{1\%}}{\overline{P}_\text{casa} - \overline{P}_\text{solar}},
\]
donde $\overline{P}$ es el promedio de los últimos 30 minutos en InfluxDB (para que una luz encendida un
momento no haga saltar la cifra), $\text{SOC}_\text{mín} = $ \code{Bateria_SocMinimo} (15\,\%) y
$E_{1\%} = $ \code{Bateria_WhPorPct} (45\,Wh, medido en las noches del 1 al 4 de noviembre de 2026
comparando la bajada del SOC con el consumo, sin sol; $\approx$\,4,5\,kWh útiles en total).

\begin{longtable}{@{}p{5.2cm}p{9.6cm}@{}}
\toprule
\textbf{Mensaje en \code{Bateria_Autonomia}} & \textbf{Condición}\\\midrule\endhead
«$\approx$ X h Y min (con N W)» & Inverter mode y consumo neto mayor de 10\,W.\\
«Cargando: el sol cubre el consumo» & Consumo neto $\le$ 10\,W.\\
«Usando la red eléctrica» & El inversor no está en modo batería.\\
«Sin datos» & Falta SOC, parámetros o lecturas.\\
\bottomrule
\end{longtable}
\begin{itemize}[nosep]
  \item \code{Bateria_Autonomia} (texto) y \code{Bateria_AutonomiaH} (horas, persistido para su gráfico).
  \item \code{Bateria_ConsumoNeto} y \code{Bateria_Disponible} (Wh hasta el mínimo).
  \item Parámetros \code{Bateria_WhPorPct} y \code{Bateria_SocMinimo}.
\end{itemize}
En la app, la autonomía se colorea en rojo por debajo de 2\,h y en naranja por debajo de 5\,h.

\section{Ahorro automático de energía}
La regla \code{ahorro-energia.js} implementa la máquina de estados de la figura~\ref{fig:ahorro}. Solo
actúa cuando el inversor funciona con batería («Inverter mode») y el Item \code{Ahorro_Auto} está
activo. \code{Ahorro_Simulacion} permite probarla sin actuar.

\begin{figure}[H]
\centering
\begin{tikzpicture}[node distance=3.3cm]
  \node[estado=verde] (n) {Normal};
  \node[estado=solar,right=of n] (b) {Bajo\\\scriptsize luces $\le$40\,\%};
  \node[estado=rojo,right=of b] (c) {Crítico\\\scriptsize cargas apagadas};
  \draw[flujo=solar] (n) to[bend left=25] node[etq,above]{SOC $\le$ 30\,\%} (b);
  \draw[flujo=rojo] (b) to[bend left=25] node[etq,above]{SOC $\le$ 20\,\%} (c);
  \draw[flujo=verde] (b) to[bend left=25] node[etq,below]{SOC $\ge$ 35\,\% o red} (n);
  \draw[flujo=verde] (c) to[bend left=40] node[etq,below]{SOC $\ge$ 35\,\% o red} (n);
  \draw[flujo=solar] (c) to[bend left=25] node[etq,below]{SOC $\ge$ 25\,\%} (b);
  \node[etq,below=1.9cm of b,text width=11cm,align=center] {Umbrales configurables desde la app: \code{Ahorro_UmbralBajo} (30), \code{Ahorro_UmbralCritico} (20), \code{Ahorro_Histeresis} (5), \code{Ahorro_NivelAtenuado} (40). La regla nunca vuelve a encender luces; cada transición envía un aviso push.};
\end{tikzpicture}
\caption{Máquina de estados del ahorro automático.}
\label{fig:ahorro}
\end{figure}

\begin{longtable}{@{}p{4.5cm}p{3.3cm}p{7cm}@{}}
\toprule
\textbf{Regla} & \textbf{Disparador} & \textbf{Acción}\\\midrule\endhead
DessMonitor: leer inversor & cron cada 2 min & Consulta la API, actualiza \code{gInversor} y \code{Energia_ApiEstado}.\\
Ahorro: evaluar batería & cambio de SOC / modo & Aplica la máquina de estados y registra \code{Ahorro_UltimaAccion}.\\
Ahorro: avisar alarmas & cambio de \code{Inversor_Alarma} & Envía un aviso push con la alarma.\\
Ahorro: vigilar conexión & periódico & Avisa si DessMonitor pasa más de 15 min sin datos.\\
\bottomrule
\end{longtable}

\section{Videovigilancia}
Todas las cámaras se leen a través del NVR (\code{.31}). Cada canal es un Thing \code{ipcamera:amcrest}
que convierte el flujo RTSP principal a MJPEG de 1280\,px y 4\,fps con \code{ffmpeg}. El servidor solo
transcodifica el MJPEG mientras alguien está mirando. Además, cada Thing genera su foto
\code{/ipcamera/<id>/ipcamera.jpg} con un \code{ffmpeg} interno (\code{snapshotUrl=ffmpeg}); la sección
«Cámaras» del sitemap muestra esas fotos refrescadas cada segundo.

\begin{figure}[H]
\centering
\begin{tikzpicture}[node distance=0.5cm and 0.7cm]
  \node[caja=lila] (cam) {Cámara\\\scriptsize HEVC/H.264};
  \node[caja=lila,right=of cam] (nvr) {NVR\\\scriptsize RTSP :554};
  \node[caja=lila,right=of nvr] (ff) {ffmpeg\\\scriptsize \code{-r 4 scale=1280}};
  \node[caja=azul,right=of ff] (ep) {\code{/ipcamera/<id>/}\\\code{ipcamera.mjpeg}};
  \node[caja=rojo,above right=0.2cm and 0.8cm of ep] (web) {MainUI\\\scriptsize url relativa};
  \node[caja=rojo,below right=0.2cm and 0.8cm of ep] (ios) {App iOS\\\scriptsize \code{/proxy?sitemap=}};
  \draw[flujo=lila] (cam) -- (nvr);
  \draw[flujo=lila] (nvr) -- (ff);
  \draw[flujo=azul] (ff) -- (ep);
  \draw[flujo=rojo] (ep) -- (web);
  \draw[flujo=rojo] (ep) -- (ios);
\end{tikzpicture}
\caption{Cadena de vídeo de las cámaras.}
\end{figure}

\begin{longtable}{@{}llll@{}}
\toprule
\textbf{Thing} & \textbf{Cámara} & \textbf{Canal NVR} & \textbf{Estado}\\\midrule\endhead
\code{ipcamera:amcrest:5a78ecb7e8} & Interior (IP8M-T2499E, PoE) & 1 & Activa\\
\code{ipcamera:amcrest:nvr_ch3} & Entrada (ASH42-B, Wi-Fi) & 3 & Activa\\
\code{ipcamera:amcrest:nvr_ch4} & Calle (IP8M-2496E, PoE) & 4 & Activa\\
\bottomrule
\end{longtable}

\paragraph{Lección aprendida (incidente del 2/11/2026).} Cuando un canal deja de emitir, cada petición
a su MJPEG bloquea un hilo de Jetty indefinidamente; las peticiones acumuladas agotaron el servidor y
toda la API REST dejó de responder. Medidas adoptadas: el vigilante que sondeaba los flujos quedó
desactivado, un canal caído se \emph{deshabilita} (responde 404 al instante) y, antes de reactivarlo, se
comprueba el RTSP desde fuera de openHAB con \code{ffprobe}.

\paragraph{Fotos congeladas (incidente del 5 al 6/11/2026).} El \code{ffmpeg} interno que genera la foto
puede colgarse: la cámara sigue \code{ONLINE} pero el binding sirve siempre la misma imagen. La regla
\code{camaras-vigilancia.js} lee cada minuto la foto local de cada cámara (solo \code{ipcamera.jpg} en el
servidor, sin abrir el vídeo de la cámara ni el NVR); si no cambia en 3 minutos, deshabilita el Thing y lo
vuelve a habilitar 5\,s después, lo que arranca un \code{ffmpeg} nuevo. Como máximo hace un reinicio por
cámara cada 15 minutos y no actúa sobre cámaras deshabilitadas o fuera de línea. El antiguo vigilante que
sondeaba los flujos MJPEG sigue desactivado (\code{camaras-vigilante.js.pausado}).

\section{Dispositivos Tuya y multimedia}
\begin{longtable}{@{}p{4.2cm}p{3.2cm}p{7.4cm}@{}}
\toprule
\textbf{Equipo} & \textbf{Binding} & \textbf{Items}\\\midrule\endhead
Interruptor Julián & tuya & \code{Interruptor_Julian_switch_1/2}\\
Interruptor ale & tuya & \code{Interruptor_Alejandro_switch_1/2} (sin conexión)\\
Panel & tuya & \code{Panel_Luz}\\
Garaje & tuya & \code{Garaje_Luz} (atenuador)\\
Sensor de gas & tuya & \code{Gas_sensor_gas_sensor_value}\\
TV de Sala / MIBOX3 / TV Alejandro & androidtv (Google TV) & \code{*_Encendido}, \code{*_Volumen}, \code{*_Silencio}, \code{*_Reproductor}, \code{*_Tecla}, \code{*_App}, \code{*_PIN}\\
TV Alejandro & chromecast & \code{TV_Ale_CC_Control}, \code{TV_Ale_CC_Volumen}, \code{TV_Ale_CC_App}, \code{TV_Ale_CC_Titulo}\\
\bottomrule
\end{longtable}

\section{Portón y luz exterior (eWeLink)}
Los equipos Sonoff se controlan a través de la nube eWeLink. La app de desarrollador de eWeLink (rol
Standard) no admite el inicio de sesión con contraseña que usa el binding \code{sonoff}, pero sí OAuth
2.0, por lo que la integración la hace la regla \code{ewelink.js}. El Thing \code{sonoff:account:ewelink}
solo guarda el App ID y el App Secret que lee la regla.

{\small\begin{longtable}{@{}p{3.3cm}p{4.4cm}p{7.1cm}@{}}
\toprule
\textbf{Equipo} & \textbf{Item} & \textbf{Funcionamiento}\\\midrule\endhead
Portón — Sonoff DUAL R2 (\code{1000ba53cf}) & \code{Ewelink_1000ba53cf_1}, \code{Ewelink_1000ba53cf_2} & Canal 1 «Abrir», canal 2 «Cerrar». Ambos en modo pulso de 1\,s configurado en eWeLink: \code{ON} da el pulso y el relé vuelve solo a \code{OFF}.\\
Luz exterior — Sonoff 4CH (\code{10000f0c77}) & \code{Ewelink_10000f0c77_2} & Solo el canal 2 tiene carga; los canales 1, 3 y 4 están libres. Incluida en el apagado general y en la persistencia.\\
\bottomrule
\end{longtable}}

\paragraph{Convención de nombres.} Cada \code{Switch} del grupo \code{gEwelinkSw} llamado
\code{Ewelink_<deviceid>} (equipo de un canal) o \code{Ewelink_<deviceid>_<n>} (canal $n$, desde 1) queda
enlazado automáticamente: basta con crear el Item en \code{ewelink.items} con \code{autoupdate="false"}.

\begin{longtable}{@{}p{4.8cm}p{3.4cm}p{6.6cm}@{}}
\toprule
\textbf{Regla} & \textbf{Disparador} & \textbf{Acción}\\\midrule\endhead
eWeLink: enlace de inicio de sesión & cada 3 min y al arrancar & Si no hay token, genera en \code{Ewelink_LoginUrl} el enlace a la página OAuth de eWeLink.\\
eWeLink: canjear código OAuth & comando a \code{Ewelink_OAuthCode} & Canjea \texttt{región|código} por el token y lo guarda.\\
eWeLink: sincronizar equipos & cada minuto (s.\,15) & Lee el estado de los equipos y actualiza los Switches; renueva el token si hace falta.\\
eWeLink: enviar comandos & comando en \code{gEwelinkSw} & Envía \code{ON}/\code{OFF} del canal a la nube.\\
\bottomrule
\end{longtable}

\begin{figure}[H]
\centering
\begin{tikzpicture}[x=2.9cm,y=-0.62cm]
  \foreach \i/\n/\c in {0/Navegador/rojo,1/Nube eWeLink/cian,2/\code{ewelink.html}/rojo,3/\code{ewelink.js}/lila} {
    \node[caja=\c,minimum width=24mm] (e\i) at (\i,0) {\n};
    \draw[linea] (\i,0.8) -- (\i,8.6);
  }
  \draw[flujo=lila] (3,1.8) -- node[etq,above]{\code{Ewelink_LoginUrl}} (0,1.8);
  \draw[flujo=rojo] (0,3) -- node[etq,above]{inicio de sesión} (1,3);
  \draw[flujo=cian,dashed] (1,4.2) -- node[etq,above]{redirección \code{?code=}} (2,4.2);
  \draw[flujo=rojo] (2,5.4) -- node[etq,above]{POST \code{Ewelink_OAuthCode}} (3,5.4);
  \draw[flujo=lila] (3,6.6) -- node[etq,above]{canje (firma HMAC)} (1,6.6);
  \draw[flujo=cian,dashed] (1,7.8) -- node[etq,above]{token (30 días)} (3,7.8);
  \node[etq,anchor=north,text width=11cm,align=center] at (1.5,8.9) {token en \code{userdata/ewelink/token.json} (fuera del repositorio); se renueva solo};
\end{tikzpicture}
\caption{Diagrama de secuencia: autorización OAuth de eWeLink.}
\end{figure}

El estado de la conexión aparece en \code{Ewelink_Estado} («Conectado: \ldots» cuando todo va bien).
La lista de equipos, sin datos sensibles, se guarda en \code{userdata/ewelink/equipos.json}.

\section{Aspiradora robot}
Roborock S5 (IP \code{.142}) integrada con el binding \code{miio}: Thing \code{miio:vacuum:roborock}. La clave
del equipo la obtiene el Thing \code{miio:cloud:xiaomi} (cuenta Xiaomi, inicio de sesión con el QR de
Xiaomi Home); el control posterior es local.
\begin{itemize}[nosep]
  \item \textbf{Control}: \code{Aspiradora_Control} (limpiar, pausa, puntual, a la base) y
        \code{Aspiradora_Potencia} (silenciosa 101, estándar 102, turbo 103, máxima 104, mopa 105).
  \item \textbf{Estado}: estado, batería, error, área y tiempo de la limpieza actual y de la última, mapa.
  \item \textbf{Mantenimiento}: vida restante del cepillo principal, cepillo lateral, filtro y sensores; señal Wi-Fi.
\end{itemize}
Items en \code{aspiradora.items}, grupo \code{gAspiradora}; icono propio \code{cleaningrobot.svg}.

\section{Salir de casa: apagado general y monitor de luces}
La regla \code{casa-general.js} mantiene una lista de 18 luces (LCN, Tuya y la luz exterior eWeLink)
con el Thing del que depende cada una.

\begin{itemize}[nosep]
  \item \textbf{Apagado general} (\code{Casa_ApagarTodo}, \code{autoupdate="false"}): el comando \code{ON}
        apaga todas las luces y los televisores que estén encendidos (solo esos, porque el comando de energía
        podría encender uno apagado). No toca la sirena ni las persianas. 20\,s después comprueba el estado
        real y escribe en \code{Casa_ApagadoEstado} «HH:MM · todo apagado» o qué luces siguen encendidas y
        cuáles no tienen conexión. El comando \code{OFF} es una simulación que solo registra qué haría.
  \item \textbf{Monitor de luces}: ante cualquier cambio de una luz, cada minuto y al arrancar, actualiza
        \code{Luces_Encendidas} (número), \code{Luces_Lista} y \code{Luces_SinConexion}. Una luz cuenta
        como sin conexión si su Thing no está \code{ONLINE} (para eWeLink, si \code{Ewelink_Estado} no
        empieza por «Conectado»).
\end{itemize}
En la app aparece en el bloque «Salir de casa» (sitemap y página \code{overview} de MainUI), con la lista
de luces encendidas que se pueden apagar una a una.

\section{Monitor de la red}
\code{helpers/monitor-red.ps1} se ejecuta fuera de openHAB, cada 5 minutos, mediante una tarea programada
que lo lanza sin ventana a través de \code{monitor-red.vbs}. Escribe en InfluxDB (bucket \code{data}):

\begin{longtable}{@{}p{3.2cm}p{11.6cm}@{}}
\toprule
\textbf{Medida} & \textbf{Contenido}\\\midrule\endhead
\code{net_ping} & 30 pings por destino: pérdida, p50, p95 y máximo.\\
\code{net_dns} & Tiempo de resolución y éxito por servidor DNS (incluido el router).\\
\code{net_servidor} & Señal y velocidad de enlace Wi-Fi del servidor, tráfico de fondo, procesos \code{ffmpeg} y cortes de la conexión con myopenHAB (Socket.IO) en \code{openhab.log}.\\
\code{net_velocidad} & Velocidad de bajada cada 30 min (Cloudflare / OVH) y de subida cada hora.\\
\bottomrule
\end{longtable}
\code{helpers/generar-dashboard-red.ps1} genera el tablero de Grafana «Red -- Casa» (\code{red-casa.json}),
que se instala copiándolo, como administrador, a la carpeta \code{dashboards-openhab} de la instalación de
Grafana (\texttt{Program Files\textbackslash GrafanaLabs\textbackslash grafana}). Las pruebas de velocidad consumen unos 7--9\,GB al día.

% =====================================================================
\chapter{Modelo de datos e interfaces}

\section{Modelo semántico}
\begin{figure}[H]
\centering
\resizebox{\linewidth}{!}{%
\begin{tikzpicture}[
  level 1/.style={sibling distance=10.8cm,level distance=1.5cm},
  level 2/.style={sibling distance=2.8cm,level distance=1.6cm},
  every node/.style={caja=azul,font=\scriptsize,minimum height=7mm},
  edge from parent/.style={draw=gris,thick}]
  \node[font=\small\bfseries] {Casa}
    child { node[caja=verde] {Planta baja}
      child { node {Entrada\\\tiny luces, garaje, sirena} }
      child { node {Cocina\\\tiny luces, gas} }
      child { node {Comedor} }
      child { node {Sala} } }
    child { node[caja=verde] {Primera planta}
      child { node {Dormitorio\\\tiny luces, panel, persiana} }
      child { node {Julián} }
      child { node {Alejandro} } };
\end{tikzpicture}}
\caption{Jerarquía del modelo semántico de openHAB.}
\end{figure}

\section{Inventario de Items}
\begin{longtable}{@{}lrp{9.5cm}@{}}
\toprule
\textbf{Tipo} & \textbf{Cantidad} & \textbf{Ejemplos}\\\midrule\endhead
Group & 19 & \code{PlantaBaja}, \code{PrimeraPlanta}, \code{gInversor}, \code{gAhorro}, \code{gAspiradora}, \code{gEwelinkSw}\\
Dimmer & 15 & salidas LCN, \code{Garaje_Luz}, volúmenes de TV\\
Switch & 21 & interruptores Tuya y eWeLink, persianas, sirena, ajustes de ahorro, \code{Casa_ApagarTodo}\\
String & 32 & \code{Inversor_Modo}, \code{Ahorro_Estado}, \code{Bateria_Autonomia}, \code{Luces_Lista}, \code{Ewelink_Estado}, \code{Aspiradora_Estado}\\
Number / Number:* & 37 & potencia, voltajes, SOC, autonomía, gas, luces encendidas, consumibles de la aspiradora\\
Player & 4 & reproductores de TV y Chromecast\\
DateTime & 2 & \code{Energia_UltimaLectura}, \code{Aspiradora_UltimaInicio}\\
Image & 1 & \code{Aspiradora_Mapa}\\\midrule
\textbf{Total} & \textbf{131} & 61 creados en MainUI (\code{jsondb}) y 70 en archivos \code{.items}\\
\bottomrule
\end{longtable}

\section{Interfaces de usuario}
\begin{itemize}[nosep]
  \item \textbf{Sitemap \code{casa}} (app iOS): Energía (flujo animado, batería, autonomía estimada con su
        explicación e historial, datos y gráficos), Ahorro automático, Planta baja (incluye la luz exterior),
        Primera planta, Televisores, Aspiradora, Portón, Seguridad (cámaras y sirena) y Salir de casa
        (apagado general y luces encendidas). Cada luz regulable tiene interruptor y deslizador.
  \item \textbf{MainUI}: páginas \code{home}, \code{overview} (con el bloque «Salir de casa»), \code{energia},
        \code{energia_graficos}, \code{seguridad} y \code{variables}.
  \item \textbf{Páginas estáticas} (\code{/static}): \code{flujo-energia.html}, \code{energia-diaria.html},
        \code{ewelink.html} (retorno del inicio de sesión OAuth).
  \item \textbf{Grafana}: tableros «Energía -- Casa» (potencias, SOC y kWh por día) y «Red -- Casa».
  \item \textbf{Acceso remoto}: myopenHAB (\code{home.myopenhab.org}); avisos push vía openHAB Cloud.
\end{itemize}

\section{Catálogo de reglas}
{\small\begin{longtable}{@{}p{3.6cm}p{5.2cm}p{6.0cm}@{}}
\toprule
\textbf{Archivo} & \textbf{Regla} & \textbf{Disparador}\\\midrule\endhead
\code{dessmonitor.js} & DessMonitor: leer inversor & cron cada 2 min\\
\code{ahorro-energia.js} & Ahorro: evaluar batería & cambio de SOC o de modo del inversor\\
 & Ahorro: avisar alarmas del inversor & cambio de \code{Inversor_Alarma}\\
 & Ahorro: vigilar conexión DessMonitor & periódico (avisa tras 15 min sin datos)\\
\code{bateria-autonomia.js} & Batería: autonomía estimada & cada 2 min, cambio de SOC, modo o parámetros, arranque\\
\code{casa-general.js} & Casa: apagado general & comando a \code{Casa_ApagarTodo}\\
 & Casa: monitor de luces & cambio de cualquier luz, cada minuto, arranque\\
\code{camaras-vigilancia.js} & Cámaras: reiniciar si la foto se congela & cada minuto (s.\,30)\\
\code{ewelink.js} & 4 reglas (enlace, canje, sincronización, comandos) & ver sección del portón\\
\bottomrule
\end{longtable}}

% =====================================================================
\chapter{Operación y mantenimiento}

\section{Respaldo}

\begin{figure}[H]
\centering
\begin{tikzpicture}[node distance=0.6cm and 0.8cm]
  \node[caja=gris] (task) {Tarea programada\\\scriptsize domingos 03:00};
  \node[caja=azul,right=of task] (sc) {\code{respaldar.ps1}};
  \node[caja=rojo,right=of sc] (red) {Eliminar secretos\\\scriptsize \code{<REDACTED>}};
  \node[caja=verde,right=of red] (git) {git commit};
  \node[nube=gris,right=of git] (gh) {GitHub\\\scriptsize privado};
  \node[caja=gris,above=0.6cm of red,font=\scriptsize] (src) {\code{conf} · \code{jsondb} · \code{helpers} · \code{Casa.lp3}};
  \draw[flujo] (task) -- (sc); \draw[flujo] (sc) -- (red); \draw[flujo] (red) -- (git); \draw[flujo] (git) -- node[etq,above]{push} (gh);
  \draw[flujo] (src) -- (red);
\end{tikzpicture}
\caption{Proceso de respaldo automático.}
\end{figure}

\begin{itemize}[nosep]
  \item Repositorio local \texttt{Documents\textbackslash casa-domotica}, remoto privado en GitHub.
  \item Se excluyen \code{users.json}, \code{influxdb-token.txt} y \code{openhab-token.txt}; las propiedades
        \code{password}, \code{localKey}, \code{accessSecret}, \code{token}\ldots{} se sustituyen por \code{<REDACTED>}.
  \item Registro de cada ejecución en \code{respaldo.log} (no se sube).
  \item Respaldo manual: \texttt{powershell -ExecutionPolicy Bypass -File scripts\textbackslash respaldar.ps1}.
\end{itemize}

\section{Restauración}
\begin{enumerate}[nosep]
  \item Instalar openHAB 5.2.3 con Java Zulu 21 y los add-ons de \code{addons.cfg}.
  \item Con el servicio detenido, copiar \code{openhab/conf} a \texttt{C:\textbackslash openhab\textbackslash conf} y
        \code{openhab/jsondb} a \texttt{C:\textbackslash openhab\textbackslash userdata\textbackslash jsondb}.
  \item Volver a introducir las credenciales marcadas \code{<REDACTED>} (cámaras, LCN, Tuya) y los tokens
        con los scripts de \code{helpers}.
  \item Iniciar openHAB y esperar a que \code{/rest/} responda de forma estable antes de cualquier otra acción.
  \item Comprobar en el registro que no aparezca «has errors, therefore ignoring it» (persistencia).
  \item eWeLink: el token vive en \code{userdata/ewelink} y no se respalda. Para obtener uno nuevo,
        abrir el enlace que publica el Item \code{Ewelink_LoginUrl} e iniciar sesión.
  \item Aspiradora: volver a iniciar sesión en la cuenta Xiaomi (QR de Xiaomi Home) si el Thing no obtiene la clave.
  \item Monitor de red: recrear la tarea programada de \code{monitor-red.vbs} (cada 5 min) y copiar el
        tablero «Red -- Casa» a Grafana.
\end{enumerate}

\section{Procedimientos ante incidentes}
\begin{longtable}{@{}p{4.2cm}p{10.6cm}@{}}
\toprule
\textbf{Síntoma} & \textbf{Procedimiento}\\\midrule\endhead
La app no carga / REST no responde & Revisar si una cámara dejó de emitir; deshabilitar su Thing; si sigue bloqueado, reiniciar el servicio \code{openHAB} (requiere administrador).\\
Una cámara sin imagen & Comprobar el canal en la app Amcrest o con \code{ffprobe}; si falla allí, es hardware/PoE/NVR. Reactivar el Thing solo cuando el RTSP funcione.\\
Luces LCN lentas o no se apagan & Revisar \code{events.log}; no reiniciar la pasarela ni abrir conexiones de prueba a la PKE (límite de licencia).\\
Tuya «Waiting for IP address» & Verificar que el equipo esté encendido y en la red; las reglas de cortafuegos para Java 21 deben permitir UDP 6666/6667.\\
Gráficos con huecos & Buscar errores de sintaxis en \code{influxdb.persist} en \code{openhab.log}.\\
Hora del bus incorrecta & Revisar el reloj de la PKE en su página web (sobre todo tras cortes de energía).\\
Cámara con la imagen fija & Normalmente la reinicia \code{camaras-vigilancia.js} en unos 3 min; buscar «la foto no cambia» en \code{openhab.log}. Si se repite, revisar el canal con \code{ffprobe}.\\
Energía «Sin datos nuevos desde las\ldots» & El colector Eybond no sube datos a DessMonitor: revisar su Wi-Fi y alimentación; openHAB no puede corregirlo.\\
Portón o luz exterior no responden & Ver \code{Ewelink_Estado}; si no dice «Conectado», el token caducó o falta: iniciar sesión con \code{Ewelink_LoginUrl}. Comprobar en la app eWeLink que el equipo esté en línea.\\
Aspiradora sin conexión & Comprobar que esté en la base y en la red (.142); revisar el Thing de la cuenta Xiaomi.\\
«Sin conexión» en Salir de casa & El Thing de esas luces no está \code{ONLINE}; su estado real es desconocido y el apagado general no puede confirmarlas.\\
\bottomrule
\end{longtable}

\section{Seguridad}
\begin{itemize}[nosep]
  \item Credenciales solo en la base de datos de openHAB y en archivos locales; nunca en el repositorio.
  \item Pendientes: cambiar las credenciales por defecto de la pasarela LCN, eliminar el archivo de
        contraseñas en texto plano del escritorio, revocar el token de API temporal, cambiar la contraseña
        de DessMonitor, desactivar el servicio SMB de la PKE si no se usa y poner la red Wi-Fi en perfil
        «Privado».
  \item Recomendado: reservas DHCP para los equipos Tuya y TV, y un SAI (UPS) para el servidor, que ha
        sufrido varios apagones bruscos.
\end{itemize}

% =====================================================================
\chapter{Historial y hoja de ruta}

\section{Cronología}
\begin{figure}[H]
\centering
\resizebox{\linewidth}{!}{%
\begin{tikzpicture}[x=1cm]
  \draw[very thick,gris,->] (0,0) -- (22.6,0);
  \foreach \x/\d/\t/\c in {
      1.0/31 oct/{Auditoría\\Fases 0--4}/azul,
      3.9/31 oct/{openHAB\\4.1 $\to$ 5.2.3}/verde,
      6.8/1 nov/{Cámaras NVR\\Tuya · flujo}/lila,
      9.7/2 nov/{Toggles · kWh/día\\incidente cámaras}/solar,
      12.6/2 nov/{TV · hora\\bus LCN}/cian,
      15.5/3 nov/{Git · GitHub\\documentación}/rojo,
      18.4/8 nov/{Autonomía · apagado\\aspiradora · cámaras}/verde,
      21.3/15 nov/{eWeLink\\datos congelados}/cian} {
    \fill[\c] (\x,0) circle (5pt);
    \node[font=\scriptsize\bfseries,above=6pt,text=\c!80!black] at (\x,0) {\d};
    \node[font=\scriptsize,below=7pt,align=center] at (\x,0) {\t};
  }
\end{tikzpicture}}
\caption{Hitos del proyecto.}
\end{figure}

\section{Problemas conocidos}
\begin{longtable}{@{}p{5cm}p{9.8cm}@{}}
\toprule
\textbf{Problema} & \textbf{Estado / acción}\\\midrule\endhead
Latencia de acuses LCN ($\approx$48\,s) el 2/11 & Desapareció tras reiniciar openHAB; vigilar pasivamente en los registros.\\
Google TV sin emparejar & Requiere el equipo encendido para mostrar el PIN.\\
Interruptor ale sin conexión & Comprobar alimentación y Wi-Fi; el Thing lo descubrirá solo.\\
Consumo estimado 0\,W con red & Limitación del inversor (no informa carga en «Mains mode»).\\
Reloj de la PKE sin NTP & Verificar tras cortes de energía.\\
Colector Eybond deja de subir datos & Ya no se publican lecturas viejas; el corte se ve como hueco y en \code{Energia_ApiEstado}.\\
Portón y luz exterior dependen de Internet & eWeLink es solo nube: sin Internet no se pueden accionar desde openHAB.\\
Consumo de datos del monitor de red & Las pruebas de velocidad gastan 7--9\,GB/día; reducir su frecuencia si el plan lo exige.\\
\bottomrule
\end{longtable}

\section{Hoja de ruta}
\begin{enumerate}[nosep]
  \item Integración con Alexa (voz, anuncios y rutinas en los Echo de la casa).
  \item Emparejar los tres Google TV y crear escenas multimedia.
  \item Reservas DHCP y perfil de red privado; SAI para el servidor.
  \item Medición directa de consumo y de la red eléctrica (medidor dedicado).
\end{enumerate}
Completado desde la versión 1.0: vigilante de cámaras pasivo (\code{camaras-vigilancia.js}, sin sondear
los flujos de vídeo).

\end{document}
